% Options for packages loaded elsewhere
\PassOptionsToPackage{unicode}{hyperref}
\PassOptionsToPackage{hyphens}{url}
\documentclass[
  ignorenonframetext,
]{beamer}
\newif\ifbibliography
\usepackage{pgfpages}
\setbeamertemplate{caption}[numbered]
\setbeamertemplate{caption label separator}{: }
\setbeamercolor{caption name}{fg=normal text.fg}
\beamertemplatenavigationsymbolsempty
% remove section numbering
\setbeamertemplate{part page}{
  \centering
  \begin{beamercolorbox}[sep=16pt,center]{part title}
    \usebeamerfont{part title}\insertpart\par
  \end{beamercolorbox}
}
\setbeamertemplate{section page}{
  \centering
  \begin{beamercolorbox}[sep=12pt,center]{section title}
    \usebeamerfont{section title}\insertsection\par
  \end{beamercolorbox}
}
\setbeamertemplate{subsection page}{
  \centering
  \begin{beamercolorbox}[sep=8pt,center]{subsection title}
    \usebeamerfont{subsection title}\insertsubsection\par
  \end{beamercolorbox}
}
% Prevent slide breaks in the middle of a paragraph
\widowpenalties 1 10000
\raggedbottom
\AtBeginPart{
  \frame{\partpage}
}
\AtBeginSection{
  \ifbibliography
  \else
    \frame{\sectionpage}
  \fi
}
\AtBeginSubsection{
  \frame{\subsectionpage}
}
\usepackage{iftex}
\ifPDFTeX
  \usepackage[T1]{fontenc}
  \usepackage[utf8]{inputenc}
  \usepackage{textcomp} % provide euro and other symbols
\else % if luatex or xetex
  \usepackage{unicode-math} % this also loads fontspec
  \defaultfontfeatures{Scale=MatchLowercase}
  \defaultfontfeatures[\rmfamily]{Ligatures=TeX,Scale=1}
\fi
\usepackage{lmodern}
\ifPDFTeX\else
  % xetex/luatex font selection
\fi
% Use upquote if available, for straight quotes in verbatim environments
\IfFileExists{upquote.sty}{\usepackage{upquote}}{}
\IfFileExists{microtype.sty}{% use microtype if available
  \usepackage[]{microtype}
  \UseMicrotypeSet[protrusion]{basicmath} % disable protrusion for tt fonts
}{}
\makeatletter
\@ifundefined{KOMAClassName}{% if non-KOMA class
  \IfFileExists{parskip.sty}{%
    \usepackage{parskip}
  }{% else
    \setlength{\parindent}{0pt}
    \setlength{\parskip}{6pt plus 2pt minus 1pt}}
}{% if KOMA class
  \KOMAoptions{parskip=half}}
\makeatother
\setlength{\emergencystretch}{3em} % prevent overfull lines
\providecommand{\tightlist}{%
  \setlength{\itemsep}{0pt}\setlength{\parskip}{0pt}}
% Auto-generated by infrastructure/rendering/_pdf_latex_helpers.py::extract_math_font_preamble.
% Minimal header passed to Pandoc via -H so Beamer slides pick up the same math font
% as the combined-PDF path. Adjust the manuscript's preamble.md to change the font globally.
\usepackage{unicode-math}
\setmathfont{latinmodern-math.otf}

% Slide layout defaults for warning-clean scientific decks.
\usepackage{etoolbox}
\IfFileExists{xurl.sty}{\usepackage{xurl}}{}
\IfFileExists{seqsplit.sty}{\usepackage{seqsplit}}{\newcommand{\seqsplit}[1]{#1}}
\protected\def\breakseq#1{\seqsplit{#1}}
\protected\def\breaktt#1{\begingroup\ttfamily\seqsplit{#1}\endgroup}
\setlength{\emergencystretch}{6em}
\tolerance=5000
\hbadness=10000
\hfuzz=1pt
\setlength{\tabcolsep}{2pt}
\AtBeginEnvironment{longtable}{\tiny\renewcommand{\arraystretch}{0.86}\setlength{\tabcolsep}{1pt}}
\AtBeginEnvironment{tabular}{\tiny\renewcommand{\arraystretch}{0.86}\setlength{\tabcolsep}{1pt}}

% Natbib and cross-reference fallbacks — slides don't load natbib
% or cleveref, but combined-PDF manuscript prose may emit these
% commands. The fallback renders citations as a bracketed key list
% and cross-references as detokenized labels so slides don't fail on
% undefined control sequences or raw underscores. \providecommand is
% a no-op if packages load later.
\providecommand{\citep}[1]{[#1]}
\providecommand{\citet}[1]{#1}
\providecommand{\citealp}[1]{#1}
\providecommand{\citeauthor}[1]{#1}
\providecommand{\citeyear}[1]{#1}
\providecommand{\cref}[1]{\texttt{\detokenize{#1}}}
\providecommand{\Cref}[1]{\texttt{\detokenize{#1}}}

\newtheorem{proposition}{Proposition}
\newtheorem{hypothesis}{Hypothesis}
\usepackage{bookmark}
\IfFileExists{xurl.sty}{\usepackage{xurl}}{} % add URL line breaks if available
\urlstyle{same}
% fallback for those not using the hyperref driver hyperxmp:
\makeatletter
\@ifundefined{xmpquote}{\newcommand{\xmpquote}[1]{#1}}{}
\makeatother
\hypersetup{
  hidelinks,
  pdfcreator={LaTeX via pandoc}}

\author{\texorpdfstring{}{}}
\date{}

\begin{document}

\section{The Active Inference Framing of the Decision
Loop}\label{sec:active-inference}

\begin{frame}[fragile,allowframebreaks]{The decision loop}
\protect\phantomsection\label{the-decision-loop}
\breaktt{src/template\_formal/agent/agent.py}'s \texttt{Agent.decide}
scores each candidate action by a closed-form quantity, formalized here
once and cited by number everywhere else in this manuscript that uses
it:

\begin{definition}[Expected free energy of a candidate action]
For a scalar Gaussian belief \(Q(o \mid \text{action}) = \mathcal{N}(\mu_q, \sigma_q^2)\) about the
outcome a candidate action would produce, and a fixed Gaussian preference \(P(o) = \mathcal{N}(\mu_p,
\sigma_p^2)\), subject to \(\sigma_q^2 > 0\) and \(\sigma_p^2 > 0\) (enforced at construction by
\breaktt{BeliefState.\_\_post\_init\_\_}, ISC-81), the \emph{expected free energy} of the action is
the sum of a KL-divergence risk term and a differential-entropy ambiguity term over \(Q\), given in
\cref{eq:expected-free-energy}. \texttt{Agent.decide} selects the candidate action minimizing this
quantity.
\end{definition}

\[
G(\text{action}) = \mathrm{KL}\!\left[Q(o \mid \text{action}) \,\|\, P(o)\right] + \mathrm{H}\!\left[Q(o \mid \text{action})\right]
\] \{\#eq:expected-free-energy\}

Both terms of @eq:expected-free-energy are ordinary, exactly-computable
closed forms --- the univariate Gaussian KL-divergence and the
univariate Gaussian differential entropy --- and
\breaktt{tests/agent/test\_agent\_free\_energy.py} checks the
implementation against a hand-derived numeric expectation (ISC-29), not
merely a ``runs without error'' smoke assertion.

A third adversarial pass (FirstPrinciples and an independent security
review, converging on the same finding without coordinating) caught a
real gap here: \breaktt{BeliefState.variance} --- the divisor in
@eq:expected-free-energy's KL term and the argument to \(\log\) in its
entropy term --- carried no validation, so
\breaktt{BeliefState(mean=0.0,\ variance=0.0)} constructed silently and
only failed three calls later with an undocumented
\breaktt{ZeroDivisionError}. \texttt{BeliefState} now validates
\texttt{variance} (finite, \texttt{\textgreater{}\ 0.0}) at construction
(ISC-81), mirroring the same construction-time-guard pattern the storage
layer's SQL identifiers and the colony harness's
\texttt{decay}/\breaktt{sensing\_noise\_std} already use.
\breaktt{tests/agent/test\_agent\_free\_energy.py} proves the fix is
load-bearing: the \texttt{ValueError} now fires at
\breaktt{BeliefState.\_\_init\_\_}, not deep inside \(G\)'s computation.
\end{frame}

\begin{frame}[allowframebreaks]{What this borrows from Friston (2005),
and what it does not}
\protect\phantomsection\label{what-this-borrows-from-friston-2005-and-what-it-does-not}
The general framing --- that adaptive behavior can be cast as
approximate minimization of a variational free energy over beliefs about
hidden or observed states --- traces to @friston2005theory, ``A Theory
of Cortical Responses.'' That paper establishes the free-energy
principle for perception; it does not itself present the risk/ambiguity
decomposition of \emph{expected} free energy used above, which was
formalized in later active-inference literature building on it. This
template borrows only the general framing --- behavior as free-energy
minimization over a belief distribution --- and states the borrowed
formula explicitly in the source docstring rather than presenting it as
a verbatim reproduction of Friston's 2005 equations. What a
research-grade active-inference implementation would additionally
require --- hierarchical generative models, policy-conditioned rollouts
over multi-step futures, precision-weighting of prediction errors, and
message-passing (variational) inference over a structured generative
model --- is out of scope for this template and is not claimed to be
present.
\end{frame}

\begin{frame}[fragile,allowframebreaks]{From individual free-energy
minimization to collective organization}
\protect\phantomsection\label{from-individual-free-energy-minimization-to-collective-organization}
The colony's convergence behavior (@sec:results-discussion; not claimed
as emergence in the stronger sense there defined) is not produced by any
single agent's free-energy calculation; it is produced by many agents
independently minimizing \(G\) against a \textbf{shared,
environment-mediated} substrate --- the pheromone field
(\breaktt{src/template\_formal/colony/pheromone.py}). This is exactly the
bridge Ehresmann and Vanbremeersch's Memory Evolutive Systems framework
(@ehresmann2007memory) is built to describe: a hierarchy in which
lower-level components (individual agents, each running its own local
free-energy minimization) interact through a shared structure to produce
higher-level, emergent organizational patterns, formalized
category-theoretically as colimits over evolving diagrams of interacting
sub-systems. This template does not implement or check Ehresmann and
Vanbremeersch's formal colimit construction --- the connection drawn
here is, like @sec:storage-functor's functorial framing, a conceptual
bridge from an individual-level computational mechanism
(expected-free-energy minimization) to a collective-level phenomenon
(stigmergic convergence), not a claim that this template's code
constitutes a formal Memory Evolutive System.
\end{frame}

\begin{frame}[fragile,allowframebreaks]{Structural boundary the decision
loop respects}
\protect\phantomsection\label{structural-boundary-the-decision-loop-respects}
\texttt{Agent.decide} and \breaktt{Agent.record\_observation} never read
or write another agent's storage file; the type system and the
\texttt{Agent} class's public surface make this structurally true rather
than merely conventionally true (ISC-30, @sec:storage-functor). The
colony coordinator that drives many agents through repeated ticks holds
references only to each agent's public interface and to the shared
\texttt{PheromoneField} \texttt{Protocol} --- never to an agent's
internal \texttt{Database} or protocol \texttt{SessionEndpoint}
(ISC-34). This keeps the ``many local free-energy minimizers
coordinating through a shared environment, with no shared internal
state'' framing honest at the code level, not just at the level of the
manuscript's prose.
\end{frame}

\end{document}
